\documentclass[12pt,a4paper]{article}

% ── Paquetes ──────────────────────────────────────────────────────────────
\usepackage[utf8]{inputenc}
\usepackage[T1]{fontenc}
\usepackage[spanish]{babel}
\usepackage{lmodern}
\usepackage{microtype}
\usepackage{geometry}
\geometry{margin=2.5cm}

\usepackage{amsmath,amssymb}
\usepackage{graphicx}
\usepackage{subcaption}
\usepackage{booktabs}
\usepackage{array}
\usepackage{tabularx}
\usepackage{xcolor}
\usepackage{hyperref}
\usepackage{float}
\usepackage{enumitem}
\usepackage{multirow}
\usepackage{listings}
\usepackage{caption}

\hypersetup{
    colorlinks=true,
    linkcolor=blue!70!black,
    urlcolor=blue!70!black,
    citecolor=green!60!black
}

\definecolor{codeblue}{RGB}{0,102,204}
\definecolor{codegray}{RGB}{128,128,128}
\definecolor{codegreen}{RGB}{0,153,0}
\definecolor{codebg}{RGB}{248,248,248}

\lstset{
    backgroundcolor=\color{codebg},
    basicstyle=\ttfamily\footnotesize,
    keywordstyle=\color{codeblue}\bfseries,
    commentstyle=\color{codegray}\itshape,
    stringstyle=\color{codegreen},
    breaklines=true,
    frame=single,
    rulecolor=\color{gray!40},
    language=Python
}

\title{
    \vspace{-1cm}
    \Large\textbf{Diseño de Redes Neuronales Confiables}\\[0.4cm]
    \large Justicia e Incertidumbre en la Concesión de Crédito\\[0.3cm]
    \normalsize Taller B4-T1 --- Home Credit Default Risk
}
\author{
    \textbf{Grupo 7} \\[0.2cm]
    \small Instituto BME /
    Daniel Gallego Sánchez , Marco Corpa Criado , Javier Fernández Guerra
}
\date{\today}

\begin{document}
\maketitle
\tableofcontents
\newpage

%% ─── 1. INTRODUCCIÓN ──────────────────────────────────────────────────────
\section{Introducción y Objetivo}
\label{sec:intro}

El presente trabajo aborda el diseño, entrenamiento y auditoría de modelos
de clasificación neuronal para la concesión de créditos. El dataset utilizado
es el \textbf{Home Credit Default Risk} (Kaggle, 307\,507 solicitudes),
cuyo objetivo es predecir si un solicitante tendrá dificultades de pago
(\texttt{TARGET=1}, $\approx$8\%) o no (\texttt{TARGET=0}, $\approx$92\%).

Los objetivos específicos del trabajo son:

\begin{enumerate}[leftmargin=*, label=\textbf{\arabic*.}]
    \item \textbf{Arquitectura customizada}: implementar capas Keras con restricciones de dominio financiero.
    \item \textbf{Aprendizaje Justo (FAIR Loss)}: diseñar una función de coste que penalice la dependencia estadística entre las predicciones y el género del solicitante.
    \item \textbf{Progresión de arquitecturas}: explorar modelos de complejidad creciente y justificar cada elección.
    \item \textbf{Incertidumbre}: cuantificar la varianza de las predicciones y analizar su relación con la calidad de los datos de entrada.
\end{enumerate}

%% ─── 2. EDA ───────────────────────────────────────────────────────────────
\section{Análisis Exploratorio de Datos}
\label{sec:eda}

\subsection{Dataset y features seleccionadas}

Se seleccionaron las 9 columnas del enunciado:
\texttt{TARGET}, \texttt{CODE\_GENDER}, \texttt{AMT\_INCOME\_TOTAL},
\texttt{AMT\_CREDIT}, \texttt{AMT\_ANNUITY}, \texttt{DAYS\_BIRTH},
\texttt{EXT\_SOURCE\_1/2/3}.

\subsection{Desbalanceo de clases}

\begin{figure}[H]
    \centering
    \includegraphics[width=0.85\textwidth]{figs/01_target_balance.pdf}
    \caption{El 91.9\% de los solicitantes pagaron a tiempo; solo el 8.1\%
             presentó dificultades de pago (razón 11.4:1).}
    \label{fig:balance}
\end{figure}

El fuerte desbalanceo tiene consecuencias directas: la métrica de evaluación
es \textbf{AUC-ROC} (no \textit{accuracy}) y se aplican \textbf{class weights}
durante el entrenamiento (clase~1 ponderada $\approx$11.4$\times$).

\subsection{Valores ausentes y máscaras}

\begin{figure}[H]
    \centering
    \includegraphics[width=0.7\textwidth]{figs/02_missing_values.pdf}
    \caption{\texttt{EXT\_SOURCE\_1} presenta el mayor porcentaje de nulos (56.4\%),
             seguido de \texttt{EXT\_SOURCE\_3} (19.8\%).}
    \label{fig:missing}
\end{figure}

Los nulos se tratan con imputación por la mediana (calculada solo en train)
y con \textbf{máscaras binarias} (\texttt{EXT\_SOURCE\_*\_MISSING}) que
señalan si el valor original estaba ausente. Estas máscaras son la materia
prima del análisis de incertidumbre (Sección~\ref{sec:uncertainty}).

\subsection{Transformación log1p}

\begin{figure}[H]
    \centering
    \includegraphics[width=\textwidth]{figs/03_log1p_transform.pdf}
    \caption{Distribuciones antes y después de $\log(1+x)$. La transformación
             elimina la cola asimétrica extrema sin pérdida de información.}
    \label{fig:log1p}
\end{figure}

\subsection{Poder predictivo}

\begin{figure}[H]
    \centering
    \includegraphics[width=0.75\textwidth]{figs/04_ext_source_kde.pdf}
    \caption{Las tres fuentes externas muestran clara separación entre clases:
             quienes impagan tienen scores más bajos.}
    \label{fig:ext_kde}
\end{figure}

\begin{figure}[H]
    \centering
    \includegraphics[width=\textwidth]{figs/05_deciles_default_rate.pdf}
    \caption{Tasa de impago por decil. \texttt{EXT\_SOURCE\_*} son monotónas
             decrecientes. \texttt{AMT\_CREDIT} y \texttt{AMT\_ANNUITY} muestran
             una \textbf{joroba}: los créditos medios tienen más impago que los
             extremos, lo que requiere al menos dos capas ocultas.}
    \label{fig:deciles}
\end{figure}

\subsection{Disparidad de género}

\begin{figure}[H]
    \centering
    \includegraphics[width=0.45\textwidth]{figs/06_gender_default_rate.pdf}
    \caption{Los hombres tienen tasa de impago del 10.1\% frente al 7.0\%
             de las mujeres. Las 4 filas con \texttt{CODE\_GENDER=XNA} fueron eliminadas.}
    \label{fig:gender}
\end{figure}

\subsection{Interacciones no lineales}

\begin{figure}[H]
    \centering
    \includegraphics[width=0.6\textwidth]{figs/07_hexbin_ext23.pdf}
    \caption{El riesgo máximo (rojo) se concentra donde \textbf{ambas} fuentes
             son bajas simultáneamente --- efecto de interacción no capturables
             con una sola capa oculta.}
    \label{fig:hexbin}
\end{figure}

%% ─── 3. PIPELINE ──────────────────────────────────────────────────────────
\section{Preprocesado y Pipeline de Datos}
\label{sec:pipeline}

\begin{table}[H]
\centering
\caption{Transformaciones aplicadas y su justificación.}
\label{tab:transforms}
\small
\begin{tabular}{@{}lll@{}}
\toprule
\textbf{Variable(s)} & \textbf{Transformación} & \textbf{Justificación} \\
\midrule
\texttt{AMT\_INCOME\_TOTAL/CREDIT/ANNUITY} & $\log(1+x)$ & Colas asimétricas \\
\texttt{DAYS\_BIRTH} & $|x|/365$ & Edad en años positiva \\
\texttt{CODE\_GENDER} & M$\to$0, F$\to$1; XNA eliminado & Binaria limpia \\
\texttt{EXT\_SOURCE\_*} & Imputación mediana (train) & 20--56\% nulos \\
\texttt{EXT\_SOURCE\_*\_MISSING} & Máscara binaria & Señal + incertidumbre \\
\texttt{DEBT\_RATIO} & AMT\_ANNUITY / AMT\_INCOME & Dominio financiero \\
Variables continuas & \texttt{StandardScaler} (train) & Gradiente estable \\
\bottomrule
\end{tabular}
\end{table}

Split estratificado \textbf{60/20/20}: train~184\,504, val~61\,501,
test~61\,502, con tasa de impago uniforme 8.07\% en los tres conjuntos.

%% ─── 4. ARQUITECTURA ─────────────────────────────────────────────────────
\section{Arquitectura del Modelo}
\label{sec:architecture}

\subsection{Justificación del tipo de red}

Se eligió \textbf{MLP denso} porque los datos son tabulares sin estructura
espacial ni temporal. RNN, CNN y Transformer son descartados por requerir
estructura secuencial o espacial que estos datos no tienen.

\subsection{Capas custom}

\subsubsection{DebtRatioLayer}

\begin{equation}
    f(x) = \sigma\bigl(\alpha \cdot (x - \theta)\bigr), \quad \alpha \geq 0
    \label{eq:debt_layer}
\end{equation}

Inicializada con $\alpha=10$ y $\theta=0.35$ (referencia regulatoria del 35\%).
Justificación: un ratio de 0.8 y uno de 1.5 son equivalentes (``imposible pagar'');
la sigmoide aplana esa diferencia. Tras el entrenamiento, $\theta$ convergió
a 0.16, revelando que el umbral real en estos datos es inferior al regulatorio.

\subsubsection{ExtSourceLayer}

\begin{equation}
    g(\mathbf{s}) = \sigma(w_1 s_1 + w_2 s_2 + w_3 s_3 + b), \quad w_i \geq 0
    \label{eq:ext_layer}
\end{equation}

La restricción $w_i \geq 0$ garantiza monotonía negativa con el impago.
Los pesos aprendidos ($w_1=0.44$, $w_2=w_3=0.61$) son \textbf{auditables}:
el modelo asigna menos peso a \texttt{EXT\_SOURCE\_1}, coherente con su mayor
tasa de nulos (56\%).

\subsection{Progresión de modelos}

\begin{table}[H]
\centering
\caption{Comparativa de modelos con 12 features (AUC en test, class weights 11.4$\times$).}
\label{tab:models}
\begin{tabular}{@{}llrrr@{}}
\toprule
\textbf{Modelo} & \textbf{Arquitectura} & \textbf{AUC val} & \textbf{AUC test} & \textbf{Paráms} \\
\midrule
M0 & Reg. Logística               & 0.7275 & 0.7335 &     13 \\
M1 & Dense 64, ReLU               & 0.7416 & 0.7449 &    897 \\
M2 & Dense 128$\to$64             & 0.7411 & 0.7447 &  9\,985 \\
M3 & M2 + Dropout(0.3, 0.2)       & 0.7424 & 0.7452 &  9\,985 \\
M4 & M3 + DebtRatioLayer          & 0.7424 & 0.7449 &  9\,988 \\
M5 & M4 + ReduceLR                & 0.7428 & 0.7449 &  9\,988 \\
\textbf{M6} & M5 + ExtSourceLayer & \textbf{0.7421} & \textbf{0.7451} & \textbf{9\,991} \\
\midrule
M3 largo & M3 + ReduceLR + 500 épocas & 0.7419 & \textbf{0.7457} & 9\,985 \\
\bottomrule
\end{tabular}
\end{table}

\paragraph{Hallazgo del entrenamiento extendido.}
Se entrenaron M2*, M3 y M6 hasta 500 épocas con
\texttt{EarlyStopping(patience=50)} y \texttt{ReduceLROnPlateau}.
M2* sin regularización presenta overfitting severo desde la época 9 (AUC
train $\approx$0.76 vs val $\approx$0.74). M3 con Dropout alcanza el mejor
AUC (0.7457). Las capas custom no mejoran el AUC bruto en ninguna
configuración, confirmando que su valor es regulatorio e interpretativo.

\begin{figure}[H]
    \centering
    \begin{subfigure}[b]{0.48\textwidth}
        \includegraphics[width=\textwidth]{figs/curves_M0.pdf}
        \caption{M0 --- Reg. Logística}
    \end{subfigure}
    \hfill
    \begin{subfigure}[b]{0.48\textwidth}
        \includegraphics[width=\textwidth]{figs/curves_M3.pdf}
        \caption{M3 --- MLP + Dropout}
    \end{subfigure}
    \\[0.4cm]
    \begin{subfigure}[b]{0.48\textwidth}
        \includegraphics[width=\textwidth]{figs/curves_M6.pdf}
        \caption{M6 --- Dual Custom}
    \end{subfigure}
    \caption{Curvas de entrenamiento representativas.}
    \label{fig:curves}
\end{figure}

\begin{figure}[H]
    \centering
    \includegraphics[width=0.65\textwidth]{figs/arch_M6_DualCustom.pdf}
    \caption{Arquitectura M6 (Dual Custom + Dropout). Las dos ramas custom
             se concatenan con la salida densa antes de la neurona final.}
    \label{fig:arch_m6}
\end{figure}

%% ─── 5. FAIR LOSS ─────────────────────────────────────────────────────────
\section{Aprendizaje Justo --- FAIR Loss}
\label{sec:fair}

\subsection{Formulación}

\begin{equation}
    \mathcal{L}_{\text{fair}} =
    \underbrace{\text{BCE}(y, \hat{y})}_{\text{clasificación}}
    + \lambda \cdot
    \underbrace{\bigl|\mathbb{E}[\hat{y} \mid G=M] -
                       \mathbb{E}[\hat{y} \mid G=F]\bigr|}_{\text{Paridad Demográfica}}
    \label{eq:fair_loss}
\end{equation}

La \textbf{Paridad Demográfica (DP)} fue elegida frente a Equalized Odds
(requiere $y_{\text{true}}$ en el cálculo, complica el grafo),
Mutual Information (no diferenciable) y correlación (solo lineal).
DP es diferenciable, solo requiere $\hat{y}$ y $G$, y es el estándar
regulatorio de la Directiva Europea de IA y la ECOA.

\subsection{Curva de Pareto multi-semilla}

Cada valor de $\lambda$ se entrenó con \textbf{3 semillas} (42, 123, 7)
para cuantificar la varianza de inicialización.

\begin{figure}[H]
    \centering
    \includegraphics[width=\textwidth]{figs/08_pareto_fair.pdf}
    \caption{Curva de Pareto con barras de error ($\pm$std entre semillas).
             La anomalía en $\lambda=1.0$ es estructural: alta sensibilidad
             a la inicialización cuando BCE y DP tienen gradientes de peso similar.}
    \label{fig:pareto}
\end{figure}

\begin{table}[H]
\centering
\caption{Resultados Pareto multi-semilla (media $\pm$ std, evaluación en test).}
\label{tab:pareto}
\begin{tabular}{@{}rrrrrrrr@{}}
\toprule
$\lambda$ & AUC & $\pm$std & DP Gap & $\pm$std & $\hat{p}(M)$ & $\hat{p}(F)$ & Reduc.\,DP \\
\midrule
0.0 & 0.7465 & 0.0002 & 0.0295 & 0.0004 & 0.0994 & 0.0699 &    0.0\% \\
0.1 & 0.7438 & 0.0004 & 0.0084 & 0.0015 & 0.0837 & 0.0753 &   71.6\% \\
0.5 & 0.7425 & 0.0005 & 0.0051 & 0.0019 & 0.0831 & 0.0780 &   82.9\% \\
\textbf{1.0} & \textbf{0.7417} & \textbf{0.0006} & \textbf{0.0028} & \textbf{0.0020} &
\textbf{0.0778} & \textbf{0.0750} & \textbf{90.5\%} \\
2.0 & 0.7396 & 0.0007 & 0.0035 & 0.0013 & 0.0768 & 0.0733 &   88.1\% \\
5.0 & 0.7351 & 0.0003 & 0.0014 & 0.0008 & 0.0732 & 0.0718 &   95.1\% \\
\bottomrule
\end{tabular}
\end{table}

\paragraph{Modelo recomendado: $\lambda=1.0$.}
Reduce el DP Gap un \textbf{90.5\%} (de 0.0295 a 0.0028) sacrificando
solo 0.005 puntos de AUC. La curva no es perfectamente monótona por
varianza estructural del optimizador SGD cuando los objetivos compiten
con gradientes de peso similar --- comportamiento conocido en la literatura
de fairness para redes neuronales.

%% ─── 6. INCERTIDUMBRE ────────────────────────────────────────────────────
\section{Incertidumbre --- MC Dropout}
\label{sec:uncertainty}

Se utiliza \textbf{Monte Carlo Dropout} (Gal \& Ghahramani, 2016):
Dropout activo en inferencia, $N=100$ pasadas por muestra:

\begin{equation}
    \hat{\sigma}^2_i = \frac{1}{N}\sum_{t=1}^{N}(\hat{y}_i^{(t)} - \bar{y}_i)^2
    \label{eq:mc_var}
\end{equation}

Elegido frente a Redes Bayesianas (rediseño completo) y Deep Ensembles
($N$ modelos $\times$ 6 lambdas = inviable). M6 ya tiene Dropout:
coste de rediseño cero.

\begin{figure}[H]
    \centering
    \includegraphics[width=\textwidth]{figs/09_uncertainty_by_class.pdf}
    \caption{M6 base (izq.): varianza similar entre clases, no reconoce
             dificultad. FAIR $\lambda=1$ (dcha.): varianza 2.5$\times$
             mayor en impagos.}
    \label{fig:unc_class}
\end{figure}

\begin{figure}[H]
    \centering
    \includegraphics[width=\textwidth]{figs/10_uncertainty_vs_missing.pdf}
    \caption{FAIR muestra varianza creciente con el número de fuentes
             externas imputadas (+79\% con las tres ausentes). M6 base: plano.}
    \label{fig:unc_missing}
\end{figure}

\begin{table}[H]
\centering
\caption{Métricas de incertidumbre ($N=100$ pasadas MC Dropout, test).}
\label{tab:uncertainty}
\begin{tabular}{@{}lrr@{}}
\toprule
\textbf{Métrica} & \textbf{M6 base} & \textbf{FAIR $\lambda=1.0$} \\
\midrule
AUC puntual                       & 0.7451 & 0.7424 \\
AUC MC-mean                       & \textbf{0.7455} & \textbf{0.7428} \\
\midrule
Var. media TARGET=0               & 0.001422 & 0.000123 \\
Var. media TARGET=1               & 0.001306 & \textbf{0.000301} \\
Ratio var(T=1)/var(T=0)           & 0.92 (\textit{peor}) & \textbf{2.45 (\textit{mejor})} \\
\midrule
Var. 0 EXT\_SOURCE ausentes       & 0.001472 & 0.000134 \\
Var. 3 EXT\_SOURCE ausentes       & 0.001414 & \textbf{0.000241 ($+79\%$)} \\
\bottomrule
\end{tabular}
\end{table}

%% ─── 7. FEATURES EXTENDIDAS ──────────────────────────────────────────────
\section{Análisis de Features Extendidas}
\label{sec:extended}

\subsection{Motivación}

El techo de 0.745 AUC con 12 features planteaba dos hipótesis:
(a) limitación de arquitectura, o (b) limitación de información.
El entrenamiento extendido (hasta 500 épocas) descartó (a).
Se analizaron las 122 columnas del dataset con \textbf{correlación de Spearman}
y \textbf{$d$ de Cohen} para seleccionar features adicionales.

\subsection{Pipeline extendido (42 features)}

Las features se seleccionaron por Spearman $>$ 0.04 o $d$ de Cohen $>$ 0.10.
Las categóricas usan \textbf{target encoding con smoothing} para variables de
alta cardinalidad (\texttt{ORGANIZATION\_TYPE}, 58 categorías) y \textbf{one-hot}
para las de baja cardinalidad ($\leq$6 categorías):

\begin{table}[H]
\centering
\caption{Features adicionales más relevantes del pipeline extendido.}
\label{tab:ext_features}
\small
\begin{tabular}{@{}llll@{}}
\toprule
\textbf{Variable} & \textbf{Encoding} & \textbf{Spearman / Cohen $d$} \\
\midrule
\texttt{NAME\_INCOME\_TYPE}        & Target encoding & std tasa = 0.159 \\
\texttt{NAME\_FAMILY\_STATUS}      & One-hot (6 cat.) & std tasa = 0.037 \\
\texttt{NAME\_EDUCATION\_TYPE}     & One-hot (5 cat.) & std tasa = 0.036 \\
\texttt{REGION\_RATING\_CLIENT\_W\_CITY} & Escalar & Spearman = 0.061 \\
\texttt{DAYS\_LAST\_PHONE\_CHANGE} & $|x|$ + Escalar & Spearman = 0.055 \\
\texttt{DAYS\_EMPLOYED}            & Outlier 365243$\to$NaN & Spearman = 0.045 \\
\texttt{OCCUPATION\_TYPE}          & Target encoding & std tasa = 0.030 \\
\texttt{ORGANIZATION\_TYPE}        & Target encoding & 58 categorías \\
\bottomrule
\end{tabular}
\end{table}

\subsection{Resultados}

\begin{figure}[H]
    \centering
    \includegraphics[width=\textwidth]{figs/12_extended_features_comparison.pdf}
    \caption{Comparativa AUC test con 12 vs 42 features. Los colores oscuros
             corresponden a los modelos extendidos.}
    \label{fig:ext_comparison}
\end{figure}

\begin{table}[H]
\centering
\caption{Impacto de ampliar las features de 12 a 42 (AUC test).}
\label{tab:ext_results}
\begin{tabular}{@{}lrrr@{}}
\toprule
\textbf{Modelo} & \textbf{12 features} & \textbf{42 features} & \textbf{$\Delta$AUC} \\
\midrule
M0 --- Reg. Logística & 0.7335 & 0.7501 & $+$0.0166 \\
M3 --- MLP + Dropout  & 0.7457 & \textbf{0.7555} & $+$0.0098 \\
M6 --- Dual Custom    & 0.7451 & 0.7550 & $+$0.0099 \\
\bottomrule
\end{tabular}
\end{table}

El salto de 12 a 42 features produce $\approx +$0.01 AUC en los MLP y
$+$0.017 en la regresión logística. Superar 0.80 con MLPs tabulares requeriría
incorporar las tablas de historial de crédito del dataset o modelos basados
en árboles (XGBoost, LightGBM), que son el estándar empírico para datos
tabulares en competiciones Kaggle.

%% ─── 8. CONCLUSIONES ─────────────────────────────────────────────────────
\section{Conclusiones}
\label{sec:conclusions}

\begin{table}[H]
\centering
\caption{Resumen comparativo del sistema completo.}
\label{tab:final}
\begin{tabular}{@{}lrrr@{}}
\toprule
\textbf{Métrica} & \textbf{M6 (12 feat.)} & \textbf{FAIR $\lambda=1.0$} & \textbf{M3 (42 feat.)} \\
\midrule
AUC-ROC (test)          & 0.7451 & 0.7417 & \textbf{0.7555} \\
DP Gap                  & 0.0295 & \textbf{0.0028} & --- \\
Reduc. DP               & ---    & \textbf{90.5\%} & --- \\
Ratio var(T=1)/var(T=0) & 0.92   & \textbf{2.48}  & --- \\
\bottomrule
\end{tabular}
\end{table}

\begin{enumerate}[leftmargin=*, label=\textbf{\arabic*.}]
    \item \textbf{El techo de AUC era de información, no de arquitectura.}
    El entrenamiento extendido confirma que 0.745 no se debe a falta de
    convergencia. Ampliar a 42 features produce $+$0.01 AUC. La señal
    adicional existe en el dataset y es explotable.

    \item \textbf{Dropout es el componente más efectivo.}
    M3 con Dropout supera a M2* sin él en todas las configuraciones.
    Las capas custom no mejoran el AUC bruto pero son auditables:
    los pesos de \texttt{ExtSourceLayer} reflejan coherentemente la calidad
    de cada fuente externa.

    \item \textbf{FAIR Loss con $\lambda=1.0$ es el óptimo de Pareto.}
    Reduce el DP Gap un 90.5\% sacrificando 0.005 AUC. La curva no es
    monótona por varianza de inicialización cuando los gradientes de
    fairness y clasificación compiten con peso similar,
    efecto cuantificado por el experimento multi-semilla.

    \item \textbf{El modelo FAIR tiene mejor calibración de incertidumbre.}
    Es 2.5$\times$ más incierto en impagos que en buenos pagadores, y
    aumenta su varianza monotónamente cuando faltan fuentes externas
    ($+$73\% con las tres imputadas). El modelo base muestra el
    comportamiento opuesto en ambas dimensiones.
\end{enumerate}

\end{document}